\SourceLeaf{397}%
\clearpage
\renewcommand{\BookRunningTitle}{Zur Theorie der komplexen Zahlen}
\pdfbookmark[0]{Zur Theorie der komplexen Zahlen}{nachlass-complexe-zahlen}
\begin{center}\Large ZUR THEORIE DER KOMPLEXEN ZAHLEN.\end{center}
\BookRule{.15\textwidth}
\pdfbookmark[1]{I. Neue Theorie der Zerlegung der Kuben}{nachlass-complexe-I}
\begin{center}[I.]\par\medskip\large NEUE THEORIE DER ZERLEGUNG DER KUBEN.\end{center}

I. Wir nehmen an, die Gleichung $x^3+y^3+z^3=0$ habe eine Lösung, nämlich $x=a,\ y=b,\ z=c$, wobei $a,b,c$ keinen gemeinsamen Teiler haben und folglich auch paarweise teilerfremd sind. Wir setzen
\[
\begin{aligned}
b+c&=\alpha\\
c+a&=\beta\\
a+b&=\gamma
\end{aligned}
\]
wobei auch $\alpha,\beta,\gamma$ notwendig paarweise teilerfremd sind. Hätten nämlich $\alpha$ und $\beta$ einen gemeinsamen Teiler, so würde dieser auch $a^3$ und $b^3$ teilen; daher müssten auch $a$ und $b$ einen gemeinsamen Teiler haben.

Wir haben nun
\[
(\beta+\gamma-\alpha)^3+(\gamma+\alpha-\beta)^3+(\alpha+\beta-\gamma)^3=0
\]
Es gilt aber die Identität
\[
\begin{aligned}
&(\beta+\gamma-\alpha)^3+(\gamma+\alpha-\beta)^3+(\alpha+\beta-\gamma)^3\\
&=(\alpha+\beta+\gamma)^3-24\alpha\beta\gamma
\end{aligned}
\]
Folglich ist
\[
(\alpha+\beta+\gamma)^3=24\alpha\beta\gamma
\]

\SourceLeaf{398}%
Sind $\alpha,\beta,\gamma$ reelle Zahlen, so ist $\alpha+\beta+\gamma$ durch $3$ teilbar, also $\alpha+\beta+\gamma)^3$ durch $27$ und folglich $\alpha\beta\gamma$ durch $9$. Daher muss eine der Zahlen $\alpha,\beta,\gamma$, zum Beispiel $\gamma$, durch $9$ teilbar sein, also auch $c^3$ und folglich $c$ durch $3$.

Sind dagegen $\alpha,\beta,\gamma$ imaginäre Zahlen, so schließen wir, dass $\alpha+\beta+\gamma$ durch $1-\varepsilon$, folglich $24\alpha\beta\gamma$ durch $(1-\varepsilon)^3$ und somit $\alpha\beta\gamma$ durch $1-\varepsilon$ teilbar sein muss. Also ist eine der Zahlen $\alpha,\beta,\gamma$ durch $1-\varepsilon$ teilbar und folglich auch eine der Zahlen $a,b,c$.

II. Allgemein gilt die Identität
\[
\begin{aligned}
&(p+q+r)^3+(p+q\varepsilon+r\varepsilon\varepsilon)^8+(p+q\varepsilon\varepsilon+r\varepsilon)^3\\
&=27pqr+3(p+q+r)(p+q\varepsilon+r\varepsilon\varepsilon)(p+q\varepsilon\varepsilon+r\varepsilon)
\end{aligned}
\]
Ist also $p+q+r=0$, so ist
\[
(p+q\varepsilon+r\varepsilon\varepsilon)^3+(p+q\varepsilon\varepsilon+r\varepsilon)^3-27pqr=0
\]

Sind hier $p,q,r,$ selbst Kuben, nämlich der Reihe nach $=a^3,b^3,c^3$, das heißt, gibt es eine Lösung der Gleichung $x^3+y^3+z^3=0$, so setzen wir
\[
\begin{aligned}
a^3+b^3\varepsilon+c^3\varepsilon\varepsilon&=a'\\
a^3+b^3\varepsilon\varepsilon+c^3\varepsilon&=b'\\
-3abc&=c'
\end{aligned}
\]
und erhalten ebenfalls $a'^3+b'^3+c'^3=0$. Aus dieser neuen Lösung lässt sich auf dieselbe Weise eine dritte ableiten und so weiter. Man sieht leicht, dass auch die dritte Lösung aus reellen Zahlen besteht, wenn dies für die erste gilt.

Weiter ist zu bemerken: Haben $a,b,c$ keinen gemeinsamen Faktor, so gilt dies auch für $a',b',c'$, abgesehen vom Faktor $1-\varepsilon$. Es ist nämlich
\[
\begin{aligned}
\frac{a'}{1-\varepsilon}&=-\varepsilon\varepsilon a^3+\varepsilon b^3=a^2-\varepsilon c^3=-b^3+\varepsilon\varepsilon c^3\\
\frac{b'}{1-\varepsilon}&=a^3-\varepsilon b^3=-\varepsilon\varepsilon a^3+\varepsilon c^3=\varepsilon\varepsilon b^3-c^3\\
\frac{c'}{1-\varepsilon}&=(\varepsilon\varepsilon-1)abc
\end{aligned}
\]
Die ersten beiden Zahlen haben also weder mit $a$ noch mit $b$ noch mit $c$ einen gemeinsamen Faktor und können auch nicht durch $1-\varepsilon$ teilbar sein, sofern nicht $a,b,c$
\SourceLeaf{399}%
alle durch $1-\varepsilon$ teilbar sind. Daher haben sie auch keinen gemeinsamen Faktor mit der dritten Zahl.

III. Aber auch der umgekehrte Weg steht offen. Wir haben gesehen, dass eine der Größen durch $1-\varepsilon$ teilbar ist; dies sei $c$. Da man für $a$ auch $a\varepsilon$ oder $a\varepsilon\varepsilon$ und ebenso für $b$ auch $b\varepsilon$ oder $b\varepsilon\varepsilon$ einsetzen kann, dürfen wir voraussetzen, dass $a$ entweder $\equiv1$ oder $\equiv-1$ ist. Wir nehmen den ersten Fall an, da im anderen Fall $b\equiv1$ gelten würde und es nur mit $a$ vertauscht werden müsste. Wir setzen also
\[
\begin{aligned}
a&=1+3\alpha\\
b&=-1+3\beta
\end{aligned}
\]
und
\[
\begin{aligned}
\frac{a\varepsilon+b\varepsilon\varepsilon}{\varepsilon-\varepsilon\varepsilon}&=1+(\varepsilon\varepsilon-\varepsilon)(\alpha\varepsilon+\beta\varepsilon\varepsilon)=A\\
\frac{a\varepsilon\varepsilon+b\varepsilon}{\varepsilon-\varepsilon\varepsilon}&=-1+(\varepsilon\varepsilon-\varepsilon)(\alpha\varepsilon\varepsilon+\beta\varepsilon)=B\\
\frac{a+b}{\varepsilon-\varepsilon\varepsilon}&=(\varepsilon\varepsilon-\varepsilon)(\alpha+\beta)=C
\end{aligned}
\]
wobei $A+B+C=0$ und $ABC=\frac{a^3+b^3}{(\varepsilon-\varepsilon\varepsilon)^3}=\left(\frac c{\varepsilon\varepsilon-\varepsilon}\right)^3$ gilt.

\Needspace{14\baselineskip}
Da hier
\[
\begin{aligned}
a&=-\varepsilon A+\varepsilon\varepsilon B\\
b&=\varepsilon\varepsilon A-\varepsilon B
\end{aligned}
\]
gilt, können $A$ und $B$ keinen gemeinsamen Faktor haben, weil dieser sonst auch ein gemeinsamer Faktor von $a$ und $b$ wäre. Wegen $A+B+C=0$ kann folglich auch $C$ weder mit $A$ noch mit $B$ einen gemeinsamen Faktor haben. Daraus folgt leicht, dass $A$ und $B$ und somit auch $C$ Kuben sind. Denn $\left(\frac c{\varepsilon\varepsilon-\varepsilon}\right)^3$ ist durch $\varepsilon-\varepsilon\varepsilon$ und folglich auch durch $(\varepsilon-\varepsilon\varepsilon)^3$ teilbar oder $\alpha+\beta$ durch $3$; daher gilt $A\equiv1$, $B\equiv-1\ (\text{mod. }3)$.

Setzen wir nun
\[
\begin{aligned}
A&=a'^3\\
B&=b'^3\\
C&=c'^3
\end{aligned}
\]
so haben wir aus der Lösung der Gleichung $x^3+y^3+z^3=0$

\SourceLeaf{400}%
\[
\begin{aligned}
x&=a\\
y&=b\\
z&=c
\end{aligned}
\]
eine andere Lösung abgeleitet,
\[
\begin{aligned}
x&=a'\\
y&=b'\\
z&=c'
\end{aligned}
\]
\[
\text{wo }a'^3b'^3c'^3=\frac{c^3}{(\varepsilon\varepsilon-\varepsilon)^3}
\]
wobei also $c'$ den Faktor $1-\varepsilon$ einmal weniger enthält als $c$. Dies ist aber widersinnig, wenn $c$ nur durch eine bestimmte Potenz von $1-\varepsilon$ teilbar ist, das heißt, wenn $c$ von $0$ verschieden ist. Denn durch Fortsetzen dieser Operationen kämen wir sonst schließlich zu einer Lösung, bei der $z$ überhaupt nicht durch $1-\varepsilon$ teilbar wäre, im Widerspruch zu (I).
\BookRule{.3\textwidth}

Für Potenzen mit dem Exponenten $5$ kann man einen ähnlichen Weg einschlagen. Ist nämlich $a^5+b^5+c^5=0$, so setzen wir $b+c=\alpha,\ c+a=\beta,\ a+b=\gamma$ und erhalten
\[
\begin{aligned}
0&=(2a)^5+(2b)^5+(2c)^5\\
&=(\beta+\gamma-\alpha)^5+(\gamma+\alpha-\beta)^5+(\alpha+\beta-\gamma)^5\\
&=(\alpha+\beta+\gamma)^5-80\alpha\beta\gamma(\alpha\alpha+\beta\beta+\gamma\gamma)
\end{aligned}
\]

\Needspace{7\baselineskip}
Es kann aber nicht $(\alpha+\beta+\gamma)^5=80\alpha\beta\gamma(\alpha\alpha+\beta\beta+\gamma\gamma)$ gelten, ohne dass eine der Zahlen $\alpha,\beta,\gamma$ durch $1-\varepsilon$ teilbar ist. Wäre nämlich keine davon teilbar, so müssten sowohl $\alpha+\beta+\gamma$ als auch $\alpha\alpha+\beta\beta+\gamma\gamma$ durch $1-\varepsilon$ teilbar sein, folglich auch $2(\alpha\alpha+\beta\beta+\gamma\gamma)+2(\alpha+\beta+\gamma)(\alpha+\beta-\gamma)=(2\alpha+\beta)^2+3\beta\beta$, was unmöglich ist.

Man kann dies auch so darstellen. Ist $a^5+b^5+c^5=0$, so ist
\[
\begin{aligned}
4(a+b+c)^5&=5(b+c)(c+a)(a+b)\\
&\qquad[(a+2b+3c)^2+3(a+c)^2-8(a+b+c)c]\\
&=5(b+c)(c+a)(a+b)\\
&\qquad[(b-c)^2+3(b+c)^2+4(a+b+c)a]
\end{aligned}
\]
\[
4(a+b+c)^5+5abc[(b-c)^2+3(b+c)^2]=5(a+b+c)\{\dotsfour\}
\]
Im Übrigen wäre der Beweis dem vorherigen sehr ähnlich.

Versucht man aber dasselbe Vorgehen bei siebenten Potenzen, so gelingt es nicht zu beweisen, dass bei einer gegebenen Lösung
\SourceLeaf{401}%
\[
a^7+b^7+c^7=0
\]
notwendig eine der Größen $a,b,c$ durch $7$ teilbar sein muss. Es folgt nämlich nur
\[
\begin{aligned}
(\alpha+\beta+\gamma)^7={}&5\alpha\beta\gamma\{3(\alpha^4+\beta^4+\gamma^4)\\
&\qquad+10(\alpha\alpha\beta\beta+\alpha\alpha\gamma\gamma+\beta\beta\gamma\gamma)\}
\end{aligned}
\]
was gelten kann, ohne dass $\alpha,\beta,\gamma$ durch $1-\varepsilon$ teilbar sind.

Hoffentlich lässt sich dies jedoch künftig aus der Natur der Determinanten und der Einheiten ableiten.
\BookRule{.3\textwidth}
